\documentclass[10pt,a4paper]{amsart}
\usepackage{amsmath,amssymb,amsthm,graphicx}
% For commutative diagrams we use Xy-pic
\usepackage[all]{xy}

% We use 2cell for 2-commutative diagrams.
\xyoption{2cell}
% Theorem environments.
%
\theoremstyle{plain}
\newtheorem{theorem}[subsection]{Theorem}
\newtheorem{proposition}[subsection]{Proposition}
\newtheorem{lemma}[subsection]{Lemma}

\theoremstyle{definition}
\newtheorem{definition}[subsection]{Definition}
\newtheorem{example}[subsection]{Example}
\newtheorem{exercise}[subsection]{Exercise}
\newtheorem{situation}[subsection]{Situation}

\theoremstyle{remark}
\newtheorem{remark}[subsection]{Remark}
\newtheorem{remarks}[subsection]{Remarks}

\numberwithin{equation}{subsection}

% Macros
%
\def\lim{\mathop{\mathrm{lim}}\nolimits}
\def\colim{\mathop{\mathrm{colim}}\nolimits}
\def\Spec{\mathop{\mathrm{Spec}}}
\def\Hom{\mathop{\mathrm{Hom}}\nolimits}
\def\Ext{\mathop{\mathrm{Ext}}\nolimits}
\def\SheafHom{\mathop{\mathcal{H}\!\mathit{om}}\nolimits}
\def\SheafExt{\mathop{\mathcal{E}\!\mathit{xt}}\nolimits}
\def\Sch{\mathit{Sch}}
\def\Mor{\mathop{\mathrm{Mor}}\nolimits}
\def\Ob{\mathop{\mathrm{Ob}}\nolimits}
\def\Sh{\mathop{\mathit{Sh}}\nolimits}
\def\NL{\mathop{N\!L}\nolimits}
\def\CH{\mathop{\mathrm{CH}}\nolimits}
\def\proetale{{pro\text{-}\acute{e}tale}}
\def\etale{{\acute{e}tale}}
\def\QCoh{\mathit{QCoh}}
\def\Ker{\mathop{\mathrm{Ker}}}
\def\Im{\mathop{\mathrm{Im}}}
\def\Coker{\mathop{\mathrm{Coker}}}
\def\Coim{\mathop{\mathrm{Coim}}}

% Boxtimes
%
\DeclareMathSymbol{\boxtimes}{\mathbin}{AMSa}{"02}

%
% Macros for moduli stacks/spaces
%
\def\QCohstack{\mathcal{QC}\!\mathit{oh}}
\def\Cohstack{\mathcal{C}\!\mathit{oh}}
\def\Spacesstack{\mathcal{S}\!\mathit{paces}}
\def\Quotfunctor{\mathrm{Quot}}
\def\Hilbfunctor{\mathrm{Hilb}}
\def\Curvesstack{\mathcal{C}\!\mathit{urves}}
\def\Polarizedstack{\mathcal{P}\!\mathit{olarized}}
\def\Complexesstack{\mathcal{C}\!\mathit{omplexes}}
% \Pic is the operator that assigns to X its picard group, usage \Pic(X)
% \Picardstack_{X/B} denotes the Picard stack of X over B
% \Picardfunctor_{X/B} denotes the Picard functor of X over B
\def\Pic{\mathop{\mathrm{Pic}}\nolimits}
\def\Picardstack{\mathcal{P}\!\mathit{ic}}
\def\Picardfunctor{\mathrm{Pic}}
\def\Deformationcategory{\mathcal{D}\!\mathit{ef}}
\usepackage{cleveref}
\crefname{originalreference}{original reference}{original references}
\Crefname{originalreference}{Original reference}{Original references}
\makeatletter
\newlabel{lemma-mayer-vietoris}{{01EB}{source}{}{sourcepointer-1}{}}
\newlabel{lemma-mayer-vietoris@cref}{{[originalreference][0][]01EB}{[1][0][]source}{}{}{}}
\makeatother
\crefname{originalreference}{original reference}{original references}
\Crefname{originalreference}{Original reference}{Original references}
\makeatletter
\newlabel{theorem-vanishing-affine-curves}{{03SC}{source}{}{sourcepointer-2}{}}
\newlabel{theorem-vanishing-affine-curves@cref}{{[originalreference][0][]03SC}{[1][0][]source}{}{}{}}
\makeatother
\title{On a curve over an algebraically closed field torsion cohomology vanishes in degrees at least three, and in degree two when the curve is affine}
\author{The Stacks Project Authors}
\date{stacks-project:etale-cohomology.tex:local-snapshot-20261011; GNU FDL 1.2 or later}
\begin{document}
\maketitle
\noindent\textit{Editorial note.} All delivered mathematical text is the author's own, copied from the pinned Stacks Project chapter file etale-cohomology.tex; no statement or proof step was rewritten, reordered or removed.\par\medskip
\noindent\textit{Editorial note.} The counted proof is the author's complete argument: the proof is written out in full: the separated case, the induction over an affine open covering, and the Mayer-Vietoris computation of the degree three group with the vanishing on the intersection.\par\medskip
\noindent\textit{Editorial note.} The label theorem-vanishing-curves names the counted statement, and the proof environment immediately following it in the source is delivered verbatim as its proof.\par\medskip
\noindent\textit{Editorial note.} The counted statement is carried by the author's own theorem environment in the chapter source.\par\medskip
\noindent\textit{Editorial note.} The statement is an author Theorem environment, delivered verbatim; the two cited inputs are kept as external pointers with their canonical Stacks tags.\par\medskip
\noindent\textit{Allowed standard prerequisites (editorial).}\begin{itemize}
\item The precise vanishing theorem for separated curves. Same chapter, the vanishing and finiteness theorem for a separated finite type scheme of dimension one over an algebraically closed field. Supplies the separated case and the vanishing on the affine intersections.
\item Mayer-Vietoris for a covering by two open subschemes. Same chapter, the Mayer-Vietoris exact sequence for an open cover. Supplies the exact sequence used to induct on the number of affine opens.
\end{itemize}\medskip
\ifdefined\hypertarget\hypertarget{sourcepointer-1}{}\fi
\noindent\textit{Source reference.} Original reference 01EB: sources/stacks-etale-cohomology/pinned-original-source.tex, source lines 7176--7176. Prerequisite pointer invoked at proof line(s) 15756; it names a standard fact of the same chapter that the counted argument uses, and every step of the counted proof built on it is retained verbatim in the delivered body.\par\medskip
\ifdefined\hypertarget\hypertarget{sourcepointer-2}{}\fi
\noindent\textit{Source reference.} Original reference 03SC: sources/stacks-etale-cohomology/pinned-original-source.tex, source lines 15656--15656. Prerequisite pointer invoked at proof line(s) 15752; it names a standard fact of the same chapter that the counted argument uses, and every step of the counted proof built on it is retained verbatim in the delivered body.\par\medskip
\begin{theorem}
\label{theorem-vanishing-curves}
Let $X$ be a finite type, dimension $1$ scheme over an
algebraically closed field $k$. Let $\mathcal{F}$ be a torsion sheaf
on $X_\etale$. Then
$$
H_\etale^q(X, \mathcal{F}) = 0, \quad \forall q \geq 3.
$$
If $X$ affine then also $H_\etale^2(X, \mathcal{F}) = 0$.
\end{theorem}

\begin{proof}
If $X$ is separated, this follows immediately from the more precise
Theorem \ref{theorem-vanishing-affine-curves}.
If $X$ is nonseparated, choose an affine open covering
$X = X_1 \cup \ldots \cup X_n$. By induction on $n$ we may assume
the vanishing holds over $U = X_1 \cup \ldots \cup X_{n - 1}$.
Then Mayer-Vietoris (Lemma \ref{lemma-mayer-vietoris}) gives
$$
H^2_\etale(U, \mathcal{F}) \oplus H^2_\etale(X_n, \mathcal{F}) \to
H^2_\etale(U \cap X_n, \mathcal{F}) \to
H^3_\etale(X, \mathcal{F}) \to 0
$$
However, since $U \cap X_n$ is an open of an affine scheme
and hence affine by our dimension assumption, the group
$H^2_\etale(U \cap X_n, \mathcal{F})$ vanishes
by Theorem \ref{theorem-vanishing-affine-curves}.
\end{proof}

\end{document}
